\chapter{MEDF Framework Design}
\label{ch:design}

The Multi-stakeholder Ethical Decision Framework (MEDF) is designed as a comprehensive, software-enabled methodology to bridge the gap between high-level ethical principles and practical AI evaluation. The framework is built upon a conceptual architecture that guides a user through a structured workflow, from initial setup to final analysis and reporting. This chapter outlines the core design principles, the main components of the framework, and the intended user workflow.

\section{Core Design Principles}
\label{sec:design_principles}

The design of MEDF is guided by four fundamental principles:

\begin{itemize}[leftmargin=2em]
  \item \textbf{Modularity and Extensibility.} The framework is designed to be modular. Governance frameworks, stakeholder profiles, and even scoring algorithms are implemented as distinct components. This allows the system to be easily extended in the future. For example, a new governance framework can be added by creating a new YAML definition file, without altering the core application logic.

  \item \textbf{Reproducibility and Auditability.} Every evaluation run through MEDF is intended to be reproducible. Given the same set of inputs (AI system scores, stakeholder weights, and framework selections), the system will always produce the same output. This is critical for auditability and for building trust in the evaluation process. All evaluation requests and results are logged, creating a complete audit trail.

  \item \textbf{Stakeholder-Centricity.} The framework explicitly acknowledges that ethical evaluation is a multi-stakeholder problem. It is designed not to find a single \enquote{correct} answer, but to surface and navigate the tensions between different value systems. The configurable weighting system is the primary mechanism for achieving this, allowing the platform to model the unique priorities of developers, regulators, affected communities, and other relevant groups.

  \item \textbf{Quantification and Decision Support.} The framework translates qualitative ethical principles into quantifiable metrics. While recognizing that not all ethical nuances can be captured in numbers, this quantitative approach provides a clear, data-driven basis for comparison and decision-making. The goal is not to replace human judgment, but to augment it with structured, analytical insights.
\end{itemize}

\section{Framework Components}
\label{sec:components}

The MEDF workflow is organized around several conceptual components:

\begin{enumerate}[leftmargin=2em]
  \item \textbf{AI System Profile.} This is the object of evaluation. It consists of a description of the AI system and, most importantly, a set of baseline scores across the six Unified Dimensions. These scores, on a 1--7 Likert scale, represent an initial assessment of the system's ethical performance, typically provided by the development team or an initial assessor.

  \item \textbf{Governance Framework Registry.} This component contains the definitions of the ethical frameworks (ALTAI, NIST RMF, MGAF). Each framework is defined in a YAML file, which specifies its dimensions and, critically, the intrinsic weight or importance of each dimension derived from its structure (e.g., the number of clauses or requirements related to that dimension). This provides a framework-specific prior for weighting.

  \item \textbf{Stakeholder Profiles.} These profiles represent the different parties interested in the AI system's evaluation. Each profile has an ID, a role (e.g., Developer, Regulator), and a set of weights that define that stakeholder's priorities across the six Unified Dimensions. These weights represent how important each dimension is to that particular stakeholder.

  \item \textbf{Scoring Engine.} This is the core analytical component. It takes the AI system's scores, the selected frameworks, and the stakeholder profiles as input. It then computes an \enquote{effective weight} for each dimension by combining the stakeholder's priorities with the framework's intrinsic emphasis (using a product pooling method). It then uses an MCDM algorithm (like TOPSIS) to calculate a final ethical score.

  \item \textbf{Conflict and Consensus Engine.} This component analyzes the results from the scoring engine. It uses statistical methods (like Spearman's rank correlation) to measure the degree of agreement or disagreement between different stakeholders. When significant conflicts are detected, it can invoke a multi-objective optimization algorithm (NSGA-II) to find a set of Pareto-optimal consensus solutions, compromise weightings that represent the best possible trade-offs between the conflicting parties.
\end{enumerate}

\section{User Workflow}
\label{sec:workflow}

The intended workflow for an analyst using the MEDF platform is as follows:

\begin{enumerate}[leftmargin=2em]
  \item \textbf{Define the Evaluation Context.} The analyst begins by selecting the AI system to be evaluated and providing its baseline dimension scores. They also select the governance frameworks against which the system should be assessed (e.g., ALTAI and NIST RMF).

  \item \textbf{Select Stakeholders.} The analyst chooses the relevant stakeholder profiles for the evaluation. They can use the default profiles (Developer, Regulator, Affected Community) or create custom profiles with specific weightings.

  \item \textbf{Run Evaluation.} The analyst executes the evaluation. The MEDF backend calculates the per-framework and overall ethical scores, showing how the system performs from the perspective of each stakeholder and in aggregate.

  \item \textbf{Analyze Conflicts.} The analyst then moves to the conflict analysis view. The platform displays the level of agreement between each pair of stakeholders, highlighting the specific dimensions where their priorities diverge most significantly.

  \item \textbf{Explore Consensus.} If conflicts are high, the analyst can run the Pareto optimization engine. The platform will generate a set of alternative \enquote{consensus} weightings that represent viable compromises. The analyst can explore these solutions to understand the nature of the trade-offs required to resolve the ethical disagreements.

  \item \textbf{Export and Report.} The analyst can export all the results, scores, conflict matrices, and consensus solutions, into a structured data format for reporting and audit purposes. The Streamlit interface provides visualizations and tables that can be directly used in reports and presentations.
\end{enumerate}

This structured workflow is illustrated in Figure~\ref{fig:stakeholder_flow}, which shows the complete stakeholder decision flow from profile definition through conflict detection to consensus generation.

\begin{figure}[H]
\centering
\begin{tikzpicture}[
  node distance=1.2cm and 1.5cm,
  block/.style={rectangle, draw, fill=blue!8, text width=10em, text centered, rounded corners, minimum height=2.8em, font=\small},
  decision/.style={diamond, draw, fill=orange!10, text width=7em, text centered, aspect=2.2, inner sep=1pt, font=\small},
  arrow/.style={-{Stealth[length=2.5mm]}, thick},
]
  \node[block] (define) {Define Stakeholder Profiles and Weights};
  \node[block, below=of define] (select) {Select Governance Frameworks};
  \node[block, below=of select] (score) {Run Scoring Engine (TOPSIS / WSM)};
  \node[block, below=of score] (conflict) {Compute Conflict Matrix (Spearman $\rho$)};
  \node[decision, below=1.4cm of conflict] (check) {Conflict $> \theta$?};
  \node[block, below left=1.4cm and 1.5cm of check] (report) {Generate Report};
  \node[block, below right=1.4cm and 1.5cm of check] (pareto) {Run Pareto Optimization (NSGA-II)};
  \node[block, below=1.4cm of pareto] (consensus) {Present Consensus Solutions};

  \draw[arrow] (define) -- (select);
  \draw[arrow] (select) -- (score);
  \draw[arrow] (score) -- (conflict);
  \draw[arrow] (conflict) -- (check);
  \draw[arrow] (check) -- node[left, font=\small] {No} (report);
  \draw[arrow] (check) -- node[right, font=\small] {Yes} (pareto);
  \draw[arrow] (pareto) -- (consensus);
  \draw[arrow] (consensus.south) -- ++(0,-0.6) -| (report.south);
\end{tikzpicture}
\caption{Stakeholder decision flow in MEDF.}
\label{fig:stakeholder_flow}
\end{figure}

This structured workflow ensures that the ethical evaluation is comprehensive, transparent, and grounded in a defensible methodology, transforming abstract ethical debates into a concrete, data-driven decision-making process.
